% SPDX-License-Identifier: MIT
% Copyright (c) 2026 cffLaTeX contributors
%
% The cffLaTeX manual. Typeset it with `l3build doc`, or with
% `lualatex cfflatex.tex` next to cfflatex.sty, cfflatex.lua, and CITATION.cff.
\documentclass{l3doc}
\usepackage{cfflatex}
\usepackage{authblk}

% l3doc strips two characters from each verbatim line, which in .dtx files
% are "% "; this file is plain LaTeX.
\fvset{gobble=0}

% The manual is its own example: its title, abstract, and author come from the
% CITATION.cff file of this package.
\includecff{CITATION.cff}
\title{\cfftitle}
% The authors, with authblk: marks for their affiliations, which authblk lists
% below them, and email footnotes for the contacts of the citation.
\cffsetauthors[style=authblk]
% hyperref cannot take the PDF author from the \author of authblk.
\hypersetup{pdfauthor={\cffauthornames}}

% The version and date of cfflatex.sty, unless the build defines others: CI
% shows the commit or the release that the manual was built from. See
% CFFLATEX_DOC_VERSION in build.lua.
\GetFileInfo{cfflatex.sty}
\providecommand\cfflatexdocversion{\fileversion}
\providecommand\cfflatexdocdate{\filedate}
\date{Version \cfflatexdocversion\ (\cfflatexdocdate)}

\begin{document}
\maketitle

\begin{abstract}
  \cffabstract
\end{abstract}

\tableofcontents

\section{Introduction}

Software is a product of research and should be cited like one~\cite{smith2016}.
The Citation File Format (CFF)~\cite{druskat2021} makes this practical by
keeping a YAML-format file with |cff| extension, typically named |CITATION.cff|,
in the software repository. The CFF file holds the metadata needed to cite the
software which an be read by both people and programs. This practice allows
automated tools to easily discover and use the citation metadata.

The CFF file and the cited manuscript are often kept and developed together in a
single repository, the latter usually written in \LaTeX{}. Developers need to
keep duplicate metadata in sync between the CFF file and the manuscript, which
can be tedious and error-prone. This complication is even more pronounced when
the software and the manuscript are developed in parallel, the former metadata
being kept in the top level of the CFF file while the latter metadata being kept
in the subsection of the CFF file such as \texttt{preferred-citation}.

\pkg{cfflatex} relieves this burden by keeping the CFF file as the single source
of metadata and importing it into the \LaTeX{} document. The metadata, such as
the title, abstract, and authors, are imported as \LaTeX{} macros that
dynamically expand whenever the CFF file changes. This ensures that the metadata
are synced across the repository, allowing for easier maintenance and automation
of the development process.

\section{Requirements}

\pkg{cfflatex} parses the CFF file in Lua using the \pkg{lt3luabridge} and
\pkg{lua-tinyyaml} packages. It needs either
\begin{itemize}
  \item Lua\LaTeX{}, which runs the Lua code directly, or
  \item pdf\LaTeX{} with unrestricted shell escape (|--shell-escape|), which
        runs the Lua code with |texlua|.
\end{itemize}
It supports \TeX{} Live 2025 and later.

\section{Usage}

This manuscript itself is an example of using \pkg{cfflatex}. Its |CITATION.cff|
is
\VerbatimInput{CITATION.cff}

The \LaTeX{} manuscript then imports and accesses the metadata by
\begin{Verbatim}
\documentclass{l3doc}
\usepackage{cfflatex}
\includecff{CITATION.cff}
\title{\cfftitle}
\author{\cffauthor{1}{given-names}~\cffauthor{1}{family-names}}

\begin{document}
\maketitle
\begin{abstract}
  \cffabstract
\end{abstract}
\end{document}
\end{Verbatim}

\subsection{Specifying the source level}

By default, \cs{includecff} reads the top level of the CFF file. The
|source=|\meta{key} option selects another citation under the top-level
\meta{key}, such as \texttt{preferred-citation}. Any other top-level key is read
in the same way. Note that CFF~1.2.0 defines \texttt{preferred-citation} as the
only such key and allows no others, so a file with another key fails CFF
validation.

\subsection{Reading CFF fields}
\label{sec:any-field}

Once the CFF file is included, \cs{cff} reads any value of the citation by its
path. The keys and the indices, counting from~1, that lead to it, are separated
by slashes. \cs{cffcount} gives the length of a list. With the |CITATION.cff| of
this manual,
\begin{center}
  \begin{tabular}{ll}
    |\cff{title}|                  & \cff{title}                  \\
    |\cff{authors/1/family-names}| & \cff{authors/1/family-names} \\
    |\cff{keywords/2}|             & \cff{keywords/2}             \\
    |\cffcount{keywords}|          & \cffcount{keywords}          \\
  \end{tabular}
\end{center}
The path is expanded first, so |\cff{authors/\cffauthorindex/email}| reads the
email address of the current author in \cs{cffforeachauthor}.

\subsection{Iterating over authors}

\cs{cffforeachauthor} runs its argument once for each author of the citation,
and \cs{cffauthorindex} gives the index of the current author. A document can
thus list the authors without knowing how many there are. The following sets the
author list of \cs{maketitle}, with \cs{and} between the authors:
\begin{Verbatim}
\author{\cffforeachauthor{%
  \ifnum\cffauthorindex>1 \and\fi
  \cffauthor{\cffauthorindex}{given-names}~%
  \cffauthor{\cffauthorindex}{family-names}}}
\end{Verbatim}

\subsection{Multiple affiliations}
\label{sec:affiliations}

CFF gives each person one \texttt{affiliation} which is a nonempty string, and
does not natively support multiple affiliations per author. \pkg{cfflatex},
allows multiple affiliations per author by separating them with semicolons in
the \texttt{affiliation} field of the author. Affiliations of each author can be
accessed individually by \cs{cffauthoraffiliation} and related commands. Refer
to Section~\ref{sec:reference} for details.

\subsection{Declaring the authors}
\label{sec:setauthors}

\cs{cffsetauthors} declares the authors of the citation for \cs{maketitle}. It
marks each author with their affiliations, numbered among the distinct
affiliations of the citation, and with a footnote for each contact of the
citation that has the same names. With the \pkg{authblk} package, which lists
the affiliations below the authors, the title page of this manual is made by
\begin{Verbatim}
\usepackage{authblk}
\includecff{CITATION.cff}
\title{\cfftitle}
\cffsetauthors[style=authblk]
\hypersetup{pdfauthor={\cffauthornames}}
\end{Verbatim}
The default style collects the authors in the \cs{author} of the standard
classes, without the affiliations. The keys of \cs{cffsetauthors} change how a
name is made, how the authors and the affiliations are declared, and the marks
and the footnotes (Section~\ref{sec:reference}).

\subsection{Text in the metadata}

Values are inserted as text. The \TeX{} special characters
\verb+\ { } % # $ & _ ~ ^+ are printed as themselves, so \LaTeX{} markup in the
CFF file is not interpreted. Line breaks inside a value become spaces, so the
paragraphs of a multi-paragraph abstract run together, and spaces at either end
are dropped. Non-ASCII characters are kept. With pdf\LaTeX{}, use a font
encoding that has them, such as \texttt{T1}.

\subsection{Recording the CFF file}

\cs{includecff} records the CFF file as an input in the |.fls| file that the
|-recorder| option writes. Build tools that read it, such as latexmk and the
LaTeX Workshop extension of Visual Studio Code, therefore rebuild the document
when the CFF file changes.

\subsection{Checking the title page}
\label{sec:check}

By default, \pkg{cfflatex} checks that \cs{maketitle} typesets the metadata of
the citation. The \cs{title} command should use \cs{cfftitle}, and the
\cs{author} command should use \cs{cffauthor} for every author of the citation.
\cs{cff} with the paths of these values counts as well. The check cannot see
metadata that is expanded before \cs{maketitle}, such as an author list built
with \cs{xdef} in the preamble, and it needs a class that keeps the title and
the authors in \cs{@title} and \cs{@author}, as the standard classes do. In such
cases, load the package with the |check=false| option:
\begin{Verbatim}
\usepackage[check=false]{cfflatex}
\end{Verbatim}

\section{Reference}
\label{sec:reference}

The package has one option, |check|, which is |true| by default
(Section~\ref{sec:check}).

\begin{function}{\includecff}
  \begin{syntax}
    \cs{includecff} \oarg{options} \marg{file}
  \end{syntax}
  Reads the citation metadata in \meta{file}. The command can be used in the
  preamble or in the document. Each call replaces the metadata read by the
  previous one.

  \begin{description}
    \item[\texttt{source=}\textmd{\meta{key}}] reads the mapping under the
          top-level \meta{key} instead. Without this option, the citation is the
          top level of the file.
  \end{description}
\end{function}

\begin{function}[EXP]{\cff, \cffcount}
  \begin{syntax}
    \cs{cff} \marg{path}
    \cs{cffcount} \marg{path}
  \end{syntax}
  Expand to the value at \meta{path} in the citation, and to the length of the
  list at \meta{path}. \meta{path} lists the keys and the indices, counting
  from~1, that lead to the value, separated by slashes, such as
  |authors/1/family-names|. A missing path gives nothing for \cs{cff} and 0 for
  \cs{cffcount}. A list or a mapping for \cs{cff}, and a value or a mapping for
  \cs{cffcount}, is an error. |\cff{title}| is \cs{cfftitle}, and
  |\cff{authors/|\meta{index}|/|\meta{field}|}| is
  \cs{cffauthor}\marg{index}\marg{field} without the check of the index.
  Section~\ref{sec:any-field} shows examples.
\end{function}

\begin{function}[EXP]{\cfftitle}
  \begin{syntax}
    \cs{cfftitle}
  \end{syntax}
  Expands to the \texttt{title} of the citation.
\end{function}

\begin{function}[EXP]{\cffabstract}
  \begin{syntax}
    \cs{cffabstract}
  \end{syntax}
  Expands to the \texttt{abstract} of the citation.
\end{function}

\begin{function}{\cffsetauthors, \cffauthornames}
  \begin{syntax}
    \cs{cffsetauthors} \oarg{keys}
    \cs{cffauthornames}
  \end{syntax}
  \cs{cffsetauthors} declares the authors of the citation for \cs{maketitle}, as
  Section~\ref{sec:setauthors} describes. Call this function once after
  \cs{includecff}. \cs{cffauthornames} expands to the names of the authors,
  separated by commas and made by the \texttt{name} template. The \meta{keys}
  are
  \begin{description}
    \item[\texttt{style=standard}] collects the authors in one \cs{author},
          separated by \cs{and}, without the affiliations. This is the default.
    \item[\texttt{style=authblk}] declares each author with \cs{author} and each
          affiliation with \cs{affil} of the \pkg{authblk} package.
    \item[\texttt{name=}\textmd{\meta{template}}] makes the name of an author.
          |#1| takes a field of the author. The default is
          |#1{given-names}~#1{family-names}|.
    \item[\texttt{affiliation-mark=}\textmd{\meta{style}}] numbers the
          affiliations with |alph| (the default), |Alph|, |arabic|, |roman|, or
          |Roman|.
    \item[\texttt{separator=}\textmd{\meta{text}}] separates the marks of an
          author. The default is a comma.
    \item[\texttt{contact=}\textmd{\meta{template}}] is the footnote of an
          author who is a contact. |#1| takes a field of the contact. The
          default is |E-mail: #1{email}|, and an empty template gives no
          footnotes.
    \item[\texttt{note=}\textmd{\marg{index}\marg{text}}] adds a footnote with
          \meta{text} to the \meta{index}-th author, such as a current
          affiliation. \meta{index} is an integer expression, such as
          |\cffauthorcount| for the last author. The key can be given several
          times. The footnotes follow those of the contacts, in order.
    \item[\texttt{author=}\textmd{\meta{template}}] declares an author, with the
          marks |#1| and the name |#2|.
    \item[\texttt{affiliation=}\textmd{\meta{template}}] declares an
          affiliation, with the mark |#1| and the name |#2|. An empty template
          gives no affiliations.
  \end{description}
\end{function}

\begin{function}[EXP]{\cffauthorcount, \cffauthor}
  \begin{syntax}
    \cs{cffauthorcount}
    \cs{cffauthor} \marg{index} \marg{field}
  \end{syntax}
  Expand to the number of authors of the citation, and to the \meta{field} of
  the \meta{index}-th author, counting from~1. \meta{field} is a key of the
  author entry in the CFF file, such as \texttt{given-names} or
  \texttt{family-names}. A field that the author does not have expands to
  nothing.
\end{function}

\begin{function}{\cffforeachauthor, \cffauthorindex}
  \begin{syntax}
    \cs{cffforeachauthor} \marg{code}
    \cs{cffauthorindex}
  \end{syntax}
  \cs{cffforeachauthor} runs \meta{code} once for each author of the citation,
  in order, and \cs{cffauthorindex} expands to the index of the current author,
  counting from~1.
\end{function}

\begin{function}[EXP]
  {
    \cffauthoraffiliationcount, \cffauthoraffiliation,
    \cffconvertaffiliationindex
  }
  \begin{syntax}
    \cs{cffauthoraffiliationcount} \marg{index}
    \cs{cffauthoraffiliation} \marg{index} \marg{k}
    \cs{cffconvertaffiliationindex} \marg{index} \marg{k}
  \end{syntax}
  Expand to the number of affiliations of the \meta{index}-th author, which is 0
  if the author has none, to the \meta{k}-th of them, and to its number among
  the distinct affiliations of the citation, which \cs{cffaffiliation} and
  \cs{cffforeachaffiliation} use.
\end{function}

\begin{function}{\cffforeachauthoraffiliation, \cffauthoraffiliationindex}
  \begin{syntax}
    \cs{cffforeachauthoraffiliation} \marg{index} \marg{code}
    \cs{cffauthoraffiliationindex}
  \end{syntax}
  \cs{cffforeachauthoraffiliation} runs \meta{code} once for each affiliation of
  the \meta{index}-th author, in order, and \cs{cffauthoraffiliationindex}
  expands to the index of the current affiliation among those of the author,
  counting from~1.
\end{function}

\begin{function}[EXP]{\cffaffiliationcount, \cffaffiliation}
  \begin{syntax}
    \cs{cffaffiliationcount}
    \cs{cffaffiliation} \marg{n}
  \end{syntax}
  Expand to the number of distinct affiliations of the citation, and to the
  \meta{n}-th of them.
\end{function}

\begin{function}{\cffforeachaffiliation, \cffaffiliationindex}
  \begin{syntax}
    \cs{cffforeachaffiliation} \marg{code}
    \cs{cffaffiliationindex}
  \end{syntax}
  \cs{cffforeachaffiliation} runs \meta{code} once for each distinct affiliation
  of the citation, in order, and \cs{cffaffiliationindex} expands to the number
  of the current affiliation.
\end{function}

\begin{function}[EXP]{\cffcontactcount, \cffcontact}
  \begin{syntax}
    \cs{cffcontactcount}
    \cs{cffcontact} \marg{index} \marg{field}
  \end{syntax}
  Expand to the number of contacts of the citation, which is 0 if it has none,
  and to the \meta{field} of the \meta{index}-th contact, counting from~1. As
  for \cs{cffauthor}, \meta{field} is a key of the contact entry, such as
  \texttt{email}, and a field that the contact does not have expands to nothing.
\end{function}

\begin{function}{\cffforeachcontact, \cffcontactindex}
  \begin{syntax}
    \cs{cffforeachcontact} \marg{code}
    \cs{cffcontactindex}
  \end{syntax}
  \cs{cffforeachcontact} runs \meta{code} once for each contact of the citation,
  in order, and \cs{cffcontactindex} expands to the index of the current
  contact, counting from~1.
\end{function}

\section{Errors and warnings}

The following errors stop compilation.
\begin{description}
  \item[CFF file `\textmd{\meta{file}}' was not found.] \meta{file} is not in
        the current directory or the \TeX{} search path.
  \item[pdfLaTeX requires unrestricted shell escape \ldots] Use Lua\LaTeX{}, or
        run pdf\LaTeX{} with |--shell-escape|.
  \item[Cannot read CFF file: \textmd{\meta{reason}}.] The file cannot be used.
        For example, it has no top-level key given by \texttt{source}, or that
        key does not hold a mapping, or the selected citation has incorrect
        structure, or an affiliation breaks the rules of
        Section~\ref{sec:affiliations}.
  \item[Install cfflatex.lua alongside cfflatex.sty.] \pkg{cfflatex} is
        incompletely installed.
\end{description}

The following errors are reported while using the commands of
Section~\ref{sec:reference}.
\begin{description}
  \item[Call \cs{includecff} first.] No CFF file has been read yet.
  \item[style=authblk needs the authblk package.] Load \pkg{authblk} before
        \cs{cffsetauthors}.
  \item[Author index `\textmd{\meta{index}}' is out of range.] The citation has
        fewer than \meta{index} authors.
  \item[Affiliation `\textmd{\meta{k}}' is out of range.] The author or the
        citation has fewer than \meta{k} affiliations.
  \item[Contact `\textmd{\meta{index}}' is out of range.] The citation has fewer
        than \meta{index} contacts.
  \item[Not inside \cs{cffforeachauthor}.] \cs{cffauthorindex} is used outside
        the loop.
  \item[Not inside \cs{cffforeachcontact}.] \cs{cffcontactindex} is used outside
        the loop.
  \item[Not a value: \textmd{\meta{path}}.] \cs{cff} is given the path of a list
        or a mapping.
  \item[Not a list: \textmd{\meta{path}}.] \cs{cffcount} is given the path of a
        value or a mapping.
  \item[Not inside an affiliation loop.] \cs{cffaffiliationindex} is used
        outside the loops over affiliations.
  \item[Not in an author affiliation loop.] The index of an affiliation of an
        author is used outside the loop over the affiliations of the author.
\end{description}

The following warnings are issued after \cs{maketitle} by the check of
Section~\ref{sec:check}.
\begin{description}
  \item[\cs{title} does not use \cs{cfftitle}.] The title page does not show the
        title of the citation.
  \item[\cs{author} does not use author(s) \textmd{\meta{indices}} of the citation.] The
        title page does not show these authors.
\end{description}

\section{License}

\pkg{cfflatex} is distributed under the MIT License. Its source code is at
\url{https://github.com/JSS95/cfflatex}, where bugs can be reported.

% Long URLs leave ragged lines rather than stretched spaces.
\raggedright
% Number the references in the order they are first cited.
\bibliographystyle{unsrt}
\bibliography{cfflatex}

\end{document}
